\input{preamble}

% OK, start here.
%
\begin{document}

\title{Groupes de Brauer}


\maketitle

\phantomsection
\label{section-phantom}

\tableofcontents

\section{Introduction}
\label{section-introduction}

\noindent
On pourra consulter les exposés de Serre au Séminaire Cartan, voir
\cite{Serre-Cartan}. Serre renvoie lui-même à
\cite{Deuring} et \cite{ANT}. Nous avons modifié certaines démonstrations ; en particulier,
nous avons utilisé un élégant argument de Rieffel pour démontrer le théorème de Wedderburn.
Il est très probable que cette modification ne constitue pas une amélioration, et nous
encourageons vivement le lecteur à lire l'exposé original de Serre.


\section{Algèbres non commutatives}
\label{section-algebras}

\noindent
Soit $k$ un corps. Dans ce chapitre, une {\it algèbre} $A$ sur $k$ est
un anneau $A$, éventuellement non commutatif, muni d'un homomorphisme d'anneaux
$k \to A$ tel que l'image de $k$ soit contenue dans le centre de $A$ et qui
envoie $1$ sur l'élément unité de $A$. Un {\it $A$-module} est un
$A$-module à droite sur lequel l'élément unité de $A$ agit comme l'identité.

\begin{definition}
\label{definition-finite}
Soit $A$ une $k$-algèbre. Nous disons que $A$ est {\it finie} si $\dim_k(A) < \infty$.
Dans ce cas, nous écrivons $[A : k] = \dim_k(A)$.
\end{definition}

\begin{definition}
\label{definition-skew-field}
Un {\it corps gauche} est un anneau éventuellement non commutatif, muni d'un élément unité
$1$, avec $1 \not = 0$, dans lequel tout élément non nul
possède un inverse multiplicatif.
\end{definition}

\noindent
Un corps gauche est une $k$-algèbre pour un certain $k$ (par exemple pour le sous-corps premier
qu'il contient). Nous utiliserons ci-dessous que tout module sur un corps gauche
est libre : une famille maximale de vecteurs linéairement indépendants constitue une
base, et une telle famille existe d'après le lemme de Zorn.

\begin{definition}
\label{definition-simple}
Soit $A$ une $k$-algèbre.
Nous disons qu'un $A$-module $M$ est {\it simple} s'il est non nul et si
ses seuls sous-$A$-modules sont $0$ et $M$.
Nous disons que $A$ est {\it simple} si elle est non nulle et si les seuls
idéaux bilatères de $A$ sont
$0$ et $A$.
\end{definition}

\begin{definition}
\label{definition-central}
Une $k$-algèbre $A$ est {\it centrale} si le centre de $A$ est l'image de
$k \to A$.
\end{definition}

\begin{definition}
\label{definition-opposite}
Étant donnée une $k$-algèbre $A$, nous notons $A^{op}$ la $k$-algèbre obtenue
en renversant l'ordre de la multiplication dans $A$. On l'appelle
l'{\it algèbre opposée}.
\end{definition}




\section{Théorème de Wedderburn}
\label{section-wedderburn}

\noindent
L'élégant argument qui suit se trouve dans un article de Rieffel, voir
\cite{Rieffel}. La démonstration pourrait difficilement être plus simple (pour reprendre
les mots du compte rendu de Carl Faith).

\begin{lemma}
\label{lemma-rieffel}
Soit $A$ un anneau éventuellement non commutatif, muni de $1$, qui ne possède aucun
idéal bilatère non trivial. Soit $M$ un idéal à droite non nul de $A$,
et considérons $M$ comme un $A$-module à droite. Alors $A$ coïncide avec le
bicommutant de $M$.
\end{lemma}

\begin{proof}
Soit $A' = \text{End}_A(M)$, de sorte que $M$ est un $A'$-module à gauche.
Posons $A'' = \text{End}_{A'}(M)$ (le bicommutant de $M$).
Considérons $M$ comme un $A''$-module à droite\footnote{Cela
signifie que, pour $a'' \in A''$ et $m \in M$, nous avons un produit
$m a'' \in M$. En particulier, la multiplication dans $A''$
est l'opposée de celle que l'on obtiendrait en écrivant les éléments de $A''$
comme des endomorphismes agissant à gauche.}.
Soit $R : A \to A''$ l'homomorphisme naturel défini par
$mR(a) = ma$. Alors $R$ est injectif, puisque $R(1) = \text{id}_M$
et que $A$ ne possède aucun idéal bilatère non trivial. Nous affirmons que $R(M)$
est un idéal à droite de $A''$. En effet, $R(m)a'' = R(ma'')$ pour $a'' \in A''$
et $m$ dans $M$, car la multiplication {\it à gauche} sur $M$ par tout élément $n$
de $M$ définit un élément de $A'$, de sorte que
$(nm)a'' = n(ma'')$ pour tout $n$ dans $M$.
Enfin, l'idéal produit $AM$ est un idéal bilatère, et donc
$A = AM$. Ainsi $R(A) = R(A)R(M)$, si bien que $R(A)$ est un idéal à droite de $A''$.
Mais $R(A)$ contient l'élément unité de $A''$, donc $R(A) = A''$.
\end{proof}

\begin{lemma}
\label{lemma-simple-module}
Soit $A$ une $k$-algèbre. Si $A$ est finie, alors
\begin{enumerate}
\item $A$ possède un module simple,
\item tout module non nul contient un sous-module simple,
\item tout module simple sur $A$ est de dimension finie sur $k$, et
\item si $M$ est un $A$-module simple, alors $\text{End}_A(M)$ est un
corps gauche.
\end{enumerate}
\end{lemma}

\begin{proof}
Bien sûr, (1) résulte de (2), puisque $A$ est un $A$-module non nul.
Pour (2), tout sous-module de dimension (finie) minimale comme espace vectoriel sur $k$
est simple. Il en existe un de dimension finie,
car un sous-module cyclique est de dimension finie. Si $M$ est simple, alors
$mA \subset M$ est un sous-module, ce qui donne (3). Tout élément non nul
de $\text{End}_A(M)$ est un isomorphisme, d'où (4).
\end{proof}

\begin{theorem}
\label{theorem-wedderburn}
\begin{slogan}
Les algèbres simples de dimension finie sur un corps sont des algèbres de matrices sur un corps gauche.
\end{slogan}
Soit $A$ une $k$-algèbre simple et finie. Alors $A$ est une algèbre de matrices sur
une $k$-algèbre finie $K$ qui est un corps gauche.
\end{theorem}

\begin{proof}
Nous pouvons choisir un sous-module simple $M \subset A$ ; alors
la $k$-algèbre $K = \text{End}_A(M)$ est un corps gauche, voir
le lemme \ref{lemma-simple-module}.
D'après le
lemme \ref{lemma-rieffel},
nous avons $A = \text{End}_K(M)$. Comme $K$ est un corps gauche et que
$M$ est de type fini (puisque $\dim_k(M) < \infty$), nous voyons que
$M$ est libre de rang fini comme $K$-module à gauche. Il en résulte immédiatement que
$A \cong \text{Mat}(n \times n, K^{op})$.
\end{proof}



\section{Lemmes sur les algèbres}
\label{section-lemmas}

\noindent
Soit $A$ une $k$-algèbre. Soit $B \subset A$ une sous-algèbre.
Le {\it centralisateur de $B$ dans $A$} est la sous-algèbre
$$
C  = \{y \in A \mid xy = yx \text{ pour tout }x \in B\}.
$$
C'est une $k$-algèbre.

\begin{lemma}
\label{lemma-centralizer}
Soient $A$, $A'$ des $k$-algèbres. Soient $B \subset A$, $B' \subset A'$ des
sous-algèbres de centralisateurs respectifs $C$, $C'$. Alors le centralisateur de
$B \otimes_k B'$ dans $A \otimes_k A'$ est $C \otimes_k C'$.
\end{lemma}

\begin{proof}
Notons $C'' \subset A \otimes_k A'$ le centralisateur de $B \otimes_k B'$.
Il est clair que $C \otimes_k C' \subset C''$. Réciproquement, tout élément
de $C''$ commute à $B \otimes 1$ et appartient donc à $C \otimes_k A'$.
De même, $C'' \subset A \otimes_k C'$. Ainsi
$C'' \subset C \otimes_k A' \cap A \otimes_k C' = C \otimes_k C'$.
\end{proof}

\begin{lemma}
\label{lemma-center-csa}
Soit $A$ une $k$-algèbre simple de dimension finie. Alors le centre $k'$ de
$A$ est une extension finie de $k$.
\end{lemma}

\begin{proof}
Écrivons $A = \text{Mat}(n \times n, K)$ pour un corps gauche $K$ de
dimension finie sur $k$ ; voir le
théorème \ref{theorem-wedderburn}.
D'après le
lemme \ref{lemma-centralizer},
le centre de $A$ est $k \otimes_k k'$, où $k' \subset K$ est le centre de
$K$. Puisque le centre d'un corps gauche est un corps, le résultat en découle.
\end{proof}

\begin{lemma}
\label{lemma-generate-two-sided-sub}
Soit $V$ un espace vectoriel sur $k$. Soit $K$ une $k$-algèbre centrale qui
est un corps gauche. Soit $W \subset V \otimes_k K$ un sous-espace vectoriel
à gauche et à droite sur $K$. Alors $W$ est engendré, comme espace vectoriel
à gauche sur $K$, par $W \cap (V \otimes 1)$.
\end{lemma}

\begin{proof}
Soit $V' \subset V$ le sous-espace vectoriel sur $k$ engendré par les
$v \in V$ tels que $v \otimes 1 \in W$. Alors $V' \otimes_k K \subset W$ et
nous avons
$$
W/(V' \otimes_k K)  \subset  (V/V') \otimes_k K.
$$
S'il existe un vecteur non nul $\overline{v} \in V/V'$ tel que
$\overline{v} \otimes 1$ appartienne à $W/(V' \otimes_k K)$,
alors $v \otimes 1 \in W$, où $v \in V$ est un relèvement de $\overline{v}$.
Cela contredit la construction de $V'$. Nous pouvons donc remplacer
$V$ par $V/V'$ et $W$ par $W/(V' \otimes_k K)$ ; il suffit alors de montrer
que $W \cap (V \otimes 1)$ est non nul si $W$ est non nul.

\medskip\noindent
Pour le voir, choisissons un élément non nul $w \in W$ qui puisse s'écrire
$w = \sum_{i = 1, \ldots, n} v_i \otimes k_i$, avec $n$ minimal.
En multipliant à droite par $k_1^{-1}$, nous pouvons supposer $k_1 = 1$.
Si $n = 1$, la conclusion est acquise puisque $v_1 \otimes 1 \in W$.
Si $n > 1$, alors, pour tout $c \in K$,
$$
c w - w c = \sum\nolimits_{i = 2, \ldots, n} v_i \otimes (c k_i - k_i c) \in W
$$
et donc $c k_i - k_i c = 0$ par minimalité de $n$.
Il s'ensuit que $k_i$ appartient au centre de $K$, qui est $k$ par
hypothèse. Ainsi $w = (v_1 + \sum k_i v_i) \otimes 1$, ce qui contredit
la minimalité de $n$.
\end{proof}

\begin{lemma}
\label{lemma-generate-two-sided-ideal}
Soit $A$ une $k$-algèbre. Soit $K$ une $k$-algèbre centrale qui est un corps
gauche. Alors tout idéal bilatère $I \subset A \otimes_k K$ est de la forme
$J \otimes_k K$ pour un idéal bilatère $J \subset A$.
En particulier, si $A$ est simple, alors $A \otimes_k K$ l'est aussi.
\end{lemma}

\begin{proof}
Posons $J = \{a \in A \mid a \otimes 1 \in I\}$. C'est un idéal bilatère
de $A$. En outre, $I = J \otimes_k K$ d'après le
lemme \ref{lemma-generate-two-sided-sub}.
\end{proof}

\begin{lemma}
\label{lemma-matrix-algebras}
Soit $R$ un anneau, éventuellement non commutatif. Soit $n \geq 1$ un entier.
Posons $R_n = \text{Mat}(n \times n, R)$.
\begin{enumerate}
\item Les foncteurs $M \mapsto M^{\oplus n}$ et
$N \mapsto Ne_{11}$ définissent des équivalences de catégories quasi inverses
$\text{Mod}_R \leftrightarrow \text{Mod}_{R_n}$.
\item Tout idéal bilatère de $R_n$ est de la forme $IR_n$ pour un certain
idéal bilatère $I$ de $R$.
\item Le centre de $R_n$ est égal au centre de $R$.
\end{enumerate}
\end{lemma}

\begin{proof}
Le point (1) est immédiat. Si $J \subset R_n$ est un idéal bilatère, alors
$J = \bigoplus_{i,j} e_{ii}Je_{jj}$, et les ensembles de coefficients de tous les
termes $e_{ii}Je_{jj}$ coïncident et constituent un idéal bilatère $I$ de $R$.
Cela prouve (2).
Le point (3) est clair.
\end{proof}

\begin{lemma}
\label{lemma-simple-module-unique}
Soit $A$ une $k$-algèbre simple de dimension finie.
\begin{enumerate}
\item Il existe, à isomorphisme près, un unique $A$-module simple $M$.
\item Tout $A$-module de dimension finie est somme directe de copies
d'un module simple.
\item Deux $A$-modules de dimension finie sont isomorphes si et
seulement s'ils ont même dimension sur $k$.
\item Si $A = \text{Mat}(n \times n, K)$, où $K$ est un corps gauche de
dimension finie sur $k$, alors $M = K^{\oplus n}$ est un $A$-module simple et
$\text{End}_A(M) = K^{op}$.
\item Si $M$ est un $A$-module simple, alors $L = \text{End}_A(M)$ est un
corps gauche de dimension finie sur $k$ qui agit à gauche sur $M$, nous avons
$A = \text{End}_L(M)$, et les centres de $A$ et de $L$ coïncident.
En outre, $[A : k] [L : k] = \dim_k(M)^2$.
\item Pour tout $A$-module non nul $N$ de dimension finie, l'algèbre
$B = \text{End}_A(N)$ est une algèbre de matrices sur le corps gauche $L$ du
point (5). De plus, $\text{End}_B(N) = A$.
\end{enumerate}
\end{lemma}

\begin{proof}
D'après le
théorème \ref{theorem-wedderburn},
nous pouvons écrire $A = \text{Mat}(n \times n, K)$ pour un corps gauche $K$
de dimension finie sur $k$. D'après le
lemme \ref{lemma-matrix-algebras},
la catégorie des modules sur $A$ est équivalente à celle des modules sur
$K$. Les points (1), (2) et (3) en résultent, puisque tout module sur $K$ est
libre. Le point (4) résulte de ce que cette équivalence transforme le $K$-module
$K$ en $M = K^{\oplus n}$. En prenant $M = K^{\oplus n}$ dans (5),
nous voyons que $L = K^{op}$. L'assertion sur le centre de $L = K^{op}$
résulte du
lemme \ref{lemma-matrix-algebras}.
L'assertion sur $\text{End}_L(M)$ résulte de la forme explicite de $M$.
La formule des dimensions est claire.
Le point (6) résulte du fait que $N$ est isomorphe à une somme directe de
copies d'un module simple.
\end{proof}

\begin{lemma}
\label{lemma-tensor-simple}
Soient $A$ et $A'$ deux $k$-algèbres simples, dont l'une est centrale et de
dimension finie sur $k$. Alors $A \otimes_k A'$ est simple.
\end{lemma}

\begin{proof}
Supposons que $A'$ soit centrale et de dimension finie sur $k$.
Écrivons $A' = \text{Mat}(n \times n, K')$ ; voir le
théorème \ref{theorem-wedderburn}.
Le centre de $K'$ est alors $k$, et nous concluons que
$A \otimes_k K'$ est simple d'après le
lemme \ref{lemma-generate-two-sided-ideal}.
Par conséquent, $A \otimes_k A' = \text{Mat}(n \times n, A \otimes_k K')$
est simple d'après le lemme \ref{lemma-matrix-algebras}.
\end{proof}

\begin{lemma}
\label{lemma-tensor-central-simple}
Le produit tensoriel de deux algèbres centrales simples sur $k$ et de
dimension finie est une algèbre centrale simple de dimension finie.
\end{lemma}

\begin{proof}
Combiner les lemmes \ref{lemma-centralizer} et \ref{lemma-tensor-simple}.
\end{proof}

\begin{lemma}
\label{lemma-base-change}
Soit $A$ une $k$-algèbre centrale simple de dimension finie.
Soit $k'/k$ une extension de corps. Alors $A' = A \otimes_k k'$ est une
$k'$-algèbre centrale simple de dimension finie.
\end{lemma}

\begin{proof}
Combiner les lemmes \ref{lemma-centralizer} et \ref{lemma-tensor-simple}.
\end{proof}

\begin{lemma}
\label{lemma-inverse}
Soit $A$ une $k$-algèbre centrale simple de dimension finie.
Alors $A \otimes_k A^{op} \cong \text{Mat}(n \times n, k)$,
où $n = [A : k]$.
\end{lemma}

\begin{proof}
D'après le lemme \ref{lemma-tensor-central-simple}, l'algèbre
$A \otimes_k A^{op}$ est simple. L'application
$$
A \otimes_k A^{op} \longrightarrow \text{End}_k(A),\quad
a \otimes a' \longmapsto (x \mapsto axa')
$$
est donc injective. Comme la source et le but de la flèche ont même dimension,
on conclut.
\end{proof}





\section{Le groupe de Brauer d'un corps}
\label{section-brauer}

\noindent
Soit $k$ un corps. Considérons deux algèbres centrales simples $A$ et $B$,
de dimension finie sur $k$. Nous disons que $A$ et $B$ sont {\it semblables}
s'il existe $n, m > 0$ tels que
$\text{Mat}(n \times n, A) \cong \text{Mat}(m \times m, B)$
comme $k$-algèbres.

\begin{lemma}
\label{lemma-similar}
Similitude.
\begin{enumerate}
\item La similitude définit une relation d'équivalence sur l'ensemble des
classes d'isomorphisme de $k$-algèbres centrales simples de dimension finie.
\item Chaque classe de similitude contient, à isomorphisme près, un unique
corps gauche central de dimension finie sur $k$.
\item Si $A = \text{Mat}(n \times n, K)$ et
$B = \text{Mat}(m \times m, K')$ pour des corps gauches centraux $K$, $K'$
de dimension finie sur $k$, alors $A$ et $B$ sont semblables si et seulement
si $K \cong K'$ comme $k$-algèbres.
\end{enumerate}
\end{lemma}

\begin{proof}
Remarquons que, d'après le théorème de Wedderburn
(théorème \ref{theorem-wedderburn}), toute algèbre centrale simple de
dimension finie s'écrit comme une algèbre de matrices sur un corps gauche
central de dimension finie. Il suffit donc de prouver la troisième assertion.
Pour cela, il suffit de montrer que, si
$A = \text{Mat}(n \times n, K) \cong \text{Mat}(m \times m, K') = B$,
alors $K \cong K'$. En effet, si $M$ est un module simple sur $A$, nous avons
$\text{End}_A(M) = K^{op}$ ; voir le
lemme \ref{lemma-simple-module-unique}.
Par conséquent, $A \cong B$ implique $K^{op} \cong (K')^{op}$, d'où le
résultat.
\end{proof}

\noindent
Étant données deux $k$-algèbres centrales simples de dimension finie $A$ et
$B$, le produit tensoriel $A \otimes_k B$ en est encore une ; voir le
lemme \ref{lemma-tensor-central-simple}.
De plus, si $A$ est semblable à $A'$, alors $A \otimes_k B$ est semblable à
$A' \otimes_k B$, car le produit tensoriel commute à la formation des
algèbres de matrices. Le produit tensoriel définit donc une loi sur les
classes d'équivalence des algèbres centrales simples de dimension finie,
manifestement associative et commutative. Enfin, le
lemme \ref{lemma-inverse}
montre que $A \otimes_k A^{op}$ est isomorphe à une algèbre de matrices,
autrement dit que $A \otimes_k A^{op}$ appartient à la classe de similitude
de $k$. Nous obtenons ainsi un groupe abélien.

\begin{definition}
\label{definition-brauer-group}
Soit $k$ un corps. Le {\it groupe de Brauer} de $k$ est le groupe abélien des
classes de similitude des $k$-algèbres centrales simples de dimension finie
défini ci-dessus. On le note $\text{Br}(k)$.
\end{definition}

\noindent
Tout homomorphisme de corps $k \to k'$ fournit un homomorphisme de groupes
$$
\text{Br}(k) \longrightarrow \text{Br}(k'),\quad
A \longmapsto A \otimes_k k'
$$
par le lemme \ref{lemma-base-change}. Autrement dit, $\text{Br}(-)$ est un
foncteur de la catégorie des corps vers celle des groupes abéliens.
Remarquons que le groupe de Brauer d'un corps est nul si et seulement si toute
extension $k \subset K$, finie et centrale, par un corps gauche est triviale.

\begin{lemma}
\label{lemma-brauer-algebraically-closed}
Le groupe de Brauer d'un corps algébriquement clos est nul.
\end{lemma}

\begin{proof}
Soit $k \subset K$ une extension, finie et centrale, par un corps gauche.
Pour tout élément $x \in K$, le sous-anneau $k[x] \subset K$ est une
sous-$k$-algèbre commutative, intègre et de dimension finie ; c'est donc un
corps, voir Algèbre, lemme \ref{algebra-lemma-integral-over-field}.
Puisque $k$ est algébriquement clos, nous en déduisons que
$k[x] = k$. Comme $x$ était arbitraire, nous concluons que $k = K$.
\end{proof}

\begin{lemma}
\label{lemma-dimension-square}
Soit $A$ une algèbre centrale simple de dimension finie sur un corps $k$.
Alors $[A : k]$ est un carré.
\end{lemma}

\begin{proof}
Cela résulte de ce que $A \otimes_k \overline{k}$ est une algèbre de matrices
sur $\overline{k}$ d'après le
lemme \ref{lemma-brauer-algebraically-closed}.
\end{proof}
\section{Skolem-Noether}
\label{section-skolem-noether}



\begin{theorem}
\label{theorem-skolem-noether}
Soit $A$ une $k$-algèbre centrale simple de dimension finie. Soit $B$ une
$k$-algèbre simple. Soient $f, g : B \to A$ deux homomorphismes de $k$-algèbres.
Alors il existe un élément inversible $x \in A$ tel que
$f(b) = xg(b)x^{-1}$ pour tout $b \in B$.
\end{theorem}

\begin{proof}
Choisissons un $A$-module simple $M$. Posons $L = \text{End}_A(M)$.
Alors $L$ est un corps gauche de centre $k$ qui agit à gauche sur $M$; voir les
lemmes \ref{lemma-simple-module} et \ref{lemma-simple-module-unique}.
Alors $M$ possède deux structures de $B \otimes_k L^{op}$-module définies par
$m \cdot_1 (b \otimes l) = lmf(b)$ et $m \cdot_2 (b \otimes l) = lmg(b)$.
La $k$-algèbre $B \otimes_k L^{op}$ est simple d'après le
lemme \ref{lemma-tensor-simple}. Comme $B$ est simple, l'existence d'un
homomorphisme de $k$-algèbres $B \to A$ implique que $B$ est de dimension finie. Ainsi
$B \otimes_k L^{op}$ est simple de dimension finie, et les deux structures de
$B \otimes_k L^{op}$-module sur $M$
sont isomorphes d'après le lemme \ref{lemma-simple-module-unique}.
Il existe donc $\varphi : M \to M$ entrelaçant ces actions.
En particulier, $\varphi$ appartient au centralisateur de $L$, ce qui implique que
$\varphi$ est la multiplication par un certain $x \in A$; voir le
lemme \ref{lemma-simple-module-unique}. En explicitant les définitions, on voit
que $x$ résout notre problème.
\end{proof}

\begin{lemma}
\label{lemma-automorphism-inner}
Soit $A$ une $k$-algèbre centrale simple de dimension finie. Tout automorphisme de la
$k$-algèbre $A$ est intérieur. En particulier, tout automorphisme de la
$k$-algèbre $\text{Mat}(n \times n, k)$
est intérieur.
\end{lemma}

\begin{proof}
Le théorème de Skolem-Noether (théorème \ref{theorem-skolem-noether})
s'applique.
\end{proof}



\section{Le théorème du centralisateur}
\label{section-centralizer}


\begin{theorem}
\label{theorem-centralizer}
Soit $A$ une algèbre centrale simple de dimension finie sur $k$, et soit
$B$ une sous-algèbre simple de $A$. Alors
\begin{enumerate}
\item le centralisateur $C$ de $B$ dans $A$ est simple,
\item $[A : k] = [B : k][C : k]$, et
\item le centralisateur de $C$ dans $A$ est $B$.
\end{enumerate}
\end{theorem}

\begin{proof}
Dans toute cette démonstration, nous utilisons librement les résultats du
lemme \ref{lemma-simple-module-unique}.
Choisissons un $A$-module simple $M$. Posons $L = \text{End}_A(M)$.
Alors $L$ est un corps gauche de centre $k$ qui agit à gauche sur $M$
et $A = \text{End}_L(M)$.
Alors $M$ est un $B \otimes_k L^{op}$-module à droite et
$C = \text{End}_{B \otimes_k L^{op}}(M)$.
Comme l'algèbre $B \otimes_k L^{op}$ est simple d'après le
lemme \ref{lemma-tensor-simple}, on voit que $C$ est simple (encore d'après le
lemme \ref{lemma-simple-module-unique}).

\medskip\noindent
Écrivons $B \otimes_k L^{op} = \text{Mat}(m \times m, K)$ pour un certain
corps gauche $K$ de dimension finie sur $k$. Alors $C = \text{Mat}(n \times n, K^{op})$
si $M$ est isomorphe à une somme directe de $n$ copies du
$B \otimes_k L^{op}$-module simple $K^{\oplus m}$ (encore d'après le lemme). Ainsi,
$\dim_k(M) = nm [K : k]$, $[B : k] [L : k] = m^2 [K : k]$,
$[C : k] = n^2 [K : k]$, et $[A : k] [L : k] = \dim_k(M)^2$ (encore d'après
le lemme). On en déduit (2).

\medskip\noindent
L'assertion (3) découle de (2) appliquée à $C \subset A$, qui montre
que $[B : k] = [C' : k]$, où $C'$ est le centralisateur de $C$ dans $A$
(ainsi que du fait évident que $B \subset C')$.
\end{proof}

\begin{lemma}
\label{lemma-when-tensor-is-equal}
Soit $A$ une algèbre centrale simple de dimension finie sur $k$, et soit
$B$ une sous-algèbre simple de $A$. Si $B$ est une
$k$-algèbre centrale, alors $A \cong B \otimes_k C$, où $C$ est le centralisateur
de $B$ dans $A$, et ce centralisateur est central simple.
\end{lemma}

\begin{proof}
Nous avons $\dim_k(A) = \dim_k(B \otimes_k C)$ d'après le
théorème \ref{theorem-centralizer}. D'après le
lemme \ref{lemma-tensor-simple},
le produit tensoriel est simple. L'application naturelle
$B \otimes_k C \to A$ est donc injective, et par conséquent un isomorphisme.
\end{proof}

\begin{lemma}
\label{lemma-self-centralizing-subfield}
Soit $A$ une algèbre centrale simple de dimension finie sur $k$.
Si $K \subset A$ est un sous-corps, les assertions suivantes sont équivalentes :
\begin{enumerate}
\item $[A : k] = [K : k]^2$,
\item $K$ est son propre centralisateur, et
\item $K$ est un sous-anneau commutatif maximal.
\end{enumerate}
\end{lemma}

\begin{proof}
Le théorème \ref{theorem-centralizer}
montre que (1) et (2) sont équivalentes.
Il est clair que (3) et (2) sont équivalentes.
\end{proof}

\begin{lemma}
\label{lemma-maximal-subfield}
\begin{slogan}
La dimension d'un corps gauche central de dimension finie est le carré de celle
de tout sous-corps commutatif maximal.
\end{slogan}
Soit $A$ un corps gauche central de dimension finie sur $k$.
Alors tout sous-corps commutatif maximal $K \subset A$ satisfait
$[A : k] = [K : k]^2$.
\end{lemma}

\begin{proof}
C'est un cas particulier du lemme \ref{lemma-self-centralizing-subfield}.
\end{proof}





\section{Corps neutralisants}
\label{section-splitting}


\begin{definition}
\label{definition-splitting}
Soit $A$ une $k$-algèbre centrale simple de dimension finie.
On dit qu'une extension $k'/k$ {\it neutralise} $A$, ou que
$k'$ est un {\it corps neutralisant} de $A$, si $A \otimes_k k'$ est
une algèbre de matrices sur $k'$.
\end{definition}

\noindent
Autrement dit, la classe de $A$ a pour image zéro
par l'application $\text{Br}(k) \to \text{Br}(k')$.

\begin{theorem}
\label{theorem-splitting}
Soit $A$ une $k$-algèbre centrale simple de dimension finie.
Soit $k'/k$ une extension de corps finie.
Les assertions suivantes sont équivalentes :
\begin{enumerate}
\item $k'$ neutralise $A$, et
\item il existe une algèbre $B$ centrale simple de dimension finie, semblable à $A$,
telle que $k' \subset B$ et $[B : k] = [k' : k]^2$.
\end{enumerate}
\end{theorem}

\begin{proof}
Supposons (2). Il suffit de montrer que $B \otimes_k k'$ est une algèbre de
matrices. Nous savons que $B \otimes_k B^{op} \cong \text{End}_k(B)$.
Comme $k'$ est le centralisateur de $k'$ dans $B^{op}$ d'après le
lemme \ref{lemma-self-centralizing-subfield},
on voit que $B \otimes_k k'$ est le centralisateur de $k \otimes k'$
dans $B \otimes_k B^{op} = \text{End}_k(B)$. Bien entendu, ce centralisateur
n'est autre que $\text{End}_{k'}(B)$, où l'on considère $B$ comme un espace vectoriel sur $k'$
au moyen du plongement $k' \to B$. D'où le résultat.

\medskip\noindent
Supposons (1). Cela signifie qu'il existe un isomorphisme
$A \otimes_k k' \cong \text{End}_{k'}(V)$ pour un certain espace vectoriel sur $k'$, noté $V$.
Soit $B$ le centralisateur de $A$ dans $\text{End}_k(V)$. Remarquons que
$k'$ est contenu dans $B$. D'après le
lemme \ref{lemma-when-tensor-is-equal},
la somme des classes de $A$ et de $B$ est nulle dans $\text{Br}(k)$.
La formule des dimensions du
théorème \ref{theorem-centralizer}
donne
$$
[B : k] [A : k] =
\dim_k(V)^2 =
[k' : k]^2 \dim_{k'}(V)^2 =
[k' : k]^2 [A : k].
$$
Ainsi, $[B : k] = [k' : k]^2$. Nous avons donc démontré le résultat pour la
classe opposée à la classe de Brauer de $A$. Cependant, $k'$ neutralise la classe de Brauer
de $A$ si et seulement s'il neutralise
la classe de Brauer de l'algèbre opposée, ce qui conclut dans tous les cas.
\end{proof}

\begin{lemma}
\label{lemma-maximal-subfield-splits}
Tout sous-corps commutatif maximal d'un corps gauche central $K$ de dimension finie sur $k$ est
un corps neutralisant de $K$.
\end{lemma}

\begin{proof}
Combiner le lemme \ref{lemma-maximal-subfield} avec le
théorème \ref{theorem-splitting}.
\end{proof}

\begin{lemma}
\label{lemma-splitting-field-degree}
Considérons un corps gauche central $K$ de dimension finie sur $k$. Soit $d^2 = [K : k]$.
Pour tout corps neutralisant fini $k'$ de $K$, le degré $[k' : k]$ est
divisible par $d$.
\end{lemma}

\begin{proof}
D'après le théorème \ref{theorem-splitting}, il existe une algèbre centrale
simple de dimension finie, notée $B$, dans la classe de Brauer de $K$, telle que
$[B : k] = [k' : k]^2$. D'après le
lemme \ref{lemma-similar},
on a $B = \text{Mat}(n \times n, K)$ pour un certain $n$.
Alors $[k' : k]^2 = n^2d^2$, d'où le résultat.
\end{proof}

\begin{proposition}
\label{proposition-separable-splitting-field}
Considérons un corps gauche central $K$ de dimension finie sur $k$.
Il existe un sous-corps commutatif maximal $k \subset k' \subset K$ qui
est séparable sur $k$.
En particulier, toute classe de Brauer possède un corps neutralisant
fini et séparable.
\end{proposition}

\begin{proof}
Comme toute classe de Brauer est représentée par un corps gauche central de dimension
finie sur $k$, la seconde assertion découle de la
première et du
lemme \ref{lemma-maximal-subfield-splits}.

\medskip\noindent
Pour démontrer la première assertion, supposons donné un sous-corps séparable
$k' \subset K$. Alors le centralisateur $K'$ de $k'$ dans $K$
a pour centre $k'$, et le problème se ramène à trouver un sous-corps commutatif
maximal de $K'$ qui soit séparable sur $k'$. Il suffit donc de montrer que, si
$k \not = K$, on peut trouver un élément $x \in K$, $x \not \in k$,
qui soit séparable sur $k$. Cette assertion est claire en caractéristique
zéro. Nous pouvons donc supposer que $k$ est de caractéristique $p > 0$. Si le
corps de base $k$ est fini, le résultat est également clair (car les
extensions de corps finis sont toujours séparables). Nous pouvons donc supposer
que $k$ est un corps infini de caractéristique positive.

\medskip\noindent
Pour obtenir une contradiction, supposons qu'aucun élément de $K \setminus k$ ne soit séparable sur $k$.
D'après la discussion de
Corps, section \ref{fields-section-algebraic},
cela signifie que le polynôme minimal de tout $x \in K$ est de la forme
$T^q - a$, où $q$ est une puissance de $p$ et $a \in k$. Comme il est
clair que tout élément de $K$ a un polynôme minimal de degré
$\leq \dim_k(K)$, on en conclut qu'il existe une puissance fixée de $p$, notée
$q$, telle que $x^q \in k$ pour tout $x \in K$.

\medskip\noindent
Considérons l'application
$$
(-)^q : K \longrightarrow K
$$
et exprimons-la dans une $k$-base $\{a_1, \ldots, a_n\}$ de $K$
avec $a_1 = 1$. Ainsi,
$$
(\sum x_i a_i)^q = \sum f_i(x_1, \ldots, x_n)a_i.
$$
Comme la multiplication dans $K$ est $k$-bilinéaire, chaque $f_i$
est un polynôme en $x_1, \ldots, x_n$ (nous omettons les détails).
Le choix de $q$ ci-dessus et le fait que $k$ soit infini montrent que
$f_i$ est identiquement nul pour $i \geq 2$. Cette annulation subsiste donc
après extension de $k$ à sa clôture algébrique $\overline{k}$. Mais
l'algèbre $K \otimes_k \overline{k}$ est une algèbre de matrices (par exemple d'après les
lemmes \ref{lemma-base-change} et \ref{lemma-brauer-algebraically-closed}),
ce qui implique qu'il existe des éléments dont la puissance $q$-ième n'est pas centrale
(par exemple $e_{11}$). C'est la contradiction cherchée.
\end{proof}

\noindent
Les résultats précédents permettent de caractériser les algèbres centrales simples de dimension
finie comme suit.

\begin{lemma}
\label{lemma-finite-central-simple-algebra}
Soit $k$ un corps. Pour une $k$-algèbre $A$, les assertions suivantes sont équivalentes :
\begin{enumerate}
\item $A$ est une $k$-algèbre centrale simple de dimension finie,
\item $A$ est un espace vectoriel de dimension finie sur $k$, $k$ est le centre de $A$,
et $A$ n'a aucun idéal bilatère non trivial,
\item il existe $d \geq 1$ tel que
$A \otimes_k \bar k \cong \text{Mat}(d \times d, \bar k)$,
\item il existe $d \geq 1$ tel que
$A \otimes_k k^{sep} \cong \text{Mat}(d \times d, k^{sep})$,
\item il existe $d \geq 1$ et une extension galoisienne finie $k'/k$
tels que
$A \otimes_k k' \cong \text{Mat}(d \times d, k')$,
\item il existe $n \geq 1$ et un corps gauche central $K$ de dimension finie
sur $k$ tels que $A \cong \text{Mat}(n \times n, K)$.
\end{enumerate}
L'entier $d$ est appelé le {\it degré} de $A$.
\end{lemma}

\begin{proof}
L'équivalence de (1) et (2) résulte des définitions; voir la
section \ref{section-algebras}.
Supposons (1). D'après la
proposition \ref{proposition-separable-splitting-field},
il existe un corps neutralisant séparable $k \subset k'$ de $A$.
Bien entendu, une clôture galoisienne de $k'/k$ est encore un corps neutralisant.
Ainsi, (1) implique (5). Il est clair que (5) $\Rightarrow$ (4)
$\Rightarrow$ (3). Supposons (3). Alors $A \otimes_k \overline{k}$
est une $\overline{k}$-algèbre centrale simple de dimension finie, par exemple d'après le
lemme \ref{lemma-matrix-algebras}.
Cela implique immédiatement que $A$ est une $k$-algèbre centrale simple de dimension finie.
Enfin, l'équivalence de (1) et (6) est le théorème de Wedderburn; voir le
théorème \ref{theorem-wedderburn}.
\end{proof}










\begin{multicols}{2}[\section{Autres chapitres}]
\noindent
Préliminaires
\begin{enumerate}
\item \hyperref[introduction-section-phantom]{Introduction}
\item \hyperref[conventions-section-phantom]{Conventions}
\item \hyperref[sets-section-phantom]{Théorie des ensembles}
\item \hyperref[categories-section-phantom]{Catégories}
\item \hyperref[topology-section-phantom]{Topologie}
\item \hyperref[sheaves-section-phantom]{Faisceaux sur les espaces}
\item \hyperref[sites-section-phantom]{Sites et faisceaux}
\item \hyperref[stacks-section-phantom]{Champs}
\item \hyperref[fields-section-phantom]{Corps}
\item \hyperref[algebra-section-phantom]{Algèbre commutative}
\item \hyperref[brauer-section-phantom]{Groupes de Brauer}
\item \hyperref[homology-section-phantom]{Algèbre homologique}
\item \hyperref[derived-section-phantom]{Catégories dérivées}
\item \hyperref[simplicial-section-phantom]{Méthodes simpliciales}
\item \hyperref[more-algebra-section-phantom]{Compléments d'algèbre}
\item \hyperref[smoothing-section-phantom]{Lissage des morphismes d'anneaux}
\item \hyperref[modules-section-phantom]{Faisceaux de modules}
\item \hyperref[sites-modules-section-phantom]{Modules sur les sites}
\item \hyperref[injectives-section-phantom]{Objets injectifs}
\item \hyperref[cohomology-section-phantom]{Cohomologie des faisceaux}
\item \hyperref[sites-cohomology-section-phantom]{Cohomologie sur les sites}
\item \hyperref[dga-section-phantom]{Algèbre différentielle graduée}
\item \hyperref[dpa-section-phantom]{Algèbre à puissances divisées}
\item \hyperref[sdga-section-phantom]{Faisceaux différentiels gradués}
\item \hyperref[hypercovering-section-phantom]{Hyperrecouvrements}
\end{enumerate}
Schémas
\begin{enumerate}
\setcounter{enumi}{25}
\item \hyperref[schemes-section-phantom]{Schémas}
\item \hyperref[constructions-section-phantom]{Constructions de schémas}
\item \hyperref[properties-section-phantom]{Propriétés des schémas}
\item \hyperref[morphisms-section-phantom]{Morphismes de schémas}
\item \hyperref[coherent-section-phantom]{Cohomologie des schémas}
\item \hyperref[divisors-section-phantom]{Diviseurs}
\item \hyperref[limits-section-phantom]{Limites de schémas}
\item \hyperref[varieties-section-phantom]{Variétés}
\item \hyperref[topologies-section-phantom]{Topologies sur les schémas}
\item \hyperref[descent-section-phantom]{Descente}
\item \hyperref[perfect-section-phantom]{Catégories dérivées de schémas}
\item \hyperref[more-morphisms-section-phantom]{Compléments sur les morphismes}
\item \hyperref[flat-section-phantom]{Compléments sur la platitude}
\item \hyperref[groupoids-section-phantom]{Schémas en groupoïdes}
\item \hyperref[more-groupoids-section-phantom]{Compléments sur les schémas en groupoïdes}
\item \hyperref[etale-section-phantom]{Morphismes étales de schémas}
\end{enumerate}
Thèmes de théorie des schémas
\begin{enumerate}
\setcounter{enumi}{41}
\item \hyperref[chow-section-phantom]{Homologie de Chow}
\item \hyperref[intersection-section-phantom]{Théorie de l'intersection}
\item \hyperref[pic-section-phantom]{Schémas de Picard des courbes}
\item \hyperref[weil-section-phantom]{Théories cohomologiques de Weil}
\item \hyperref[adequate-section-phantom]{Modules adéquats}
\item \hyperref[dualizing-section-phantom]{Complexes dualisants}
\item \hyperref[duality-section-phantom]{Dualité pour les schémas}
\item \hyperref[discriminant-section-phantom]{Discriminants et différentes}
\item \hyperref[derham-section-phantom]{Cohomologie de de Rham}
\item \hyperref[local-cohomology-section-phantom]{Cohomologie locale}
\item \hyperref[algebraization-section-phantom]{Géométrie algébrique et formelle}
\item \hyperref[curves-section-phantom]{Courbes algébriques}
\item \hyperref[resolve-section-phantom]{Résolution des surfaces}
\item \hyperref[models-section-phantom]{Réduction semi-stable}
\item \hyperref[functors-section-phantom]{Foncteurs et morphismes}
\item \hyperref[equiv-section-phantom]{Catégories dérivées de variétés}
\item \hyperref[pione-section-phantom]{Groupes fondamentaux de schémas}
\item \hyperref[etale-cohomology-section-phantom]{Cohomologie étale}
\item \hyperref[crystalline-section-phantom]{Cohomologie cristalline}
\item \hyperref[proetale-section-phantom]{Cohomologie pro-étale}
\item \hyperref[relative-cycles-section-phantom]{Cycles relatifs}
\item \hyperref[more-etale-section-phantom]{Compléments de cohomologie étale}
\item \hyperref[trace-section-phantom]{La formule des traces}
\end{enumerate}
Espaces algébriques
\begin{enumerate}
\setcounter{enumi}{64}
\item \hyperref[spaces-section-phantom]{Espaces algébriques}
\item \hyperref[spaces-properties-section-phantom]{Propriétés des espaces algébriques}
\item \hyperref[spaces-morphisms-section-phantom]{Morphismes d'espaces algébriques}
\item \hyperref[decent-spaces-section-phantom]{Espaces algébriques décents}
\item \hyperref[spaces-cohomology-section-phantom]{Cohomologie des espaces algébriques}
\item \hyperref[spaces-limits-section-phantom]{Limites d'espaces algébriques}
\item \hyperref[spaces-divisors-section-phantom]{Diviseurs sur les espaces algébriques}
\item \hyperref[spaces-over-fields-section-phantom]{Espaces algébriques sur des corps}
\item \hyperref[spaces-topologies-section-phantom]{Topologies sur les espaces algébriques}
\item \hyperref[spaces-descent-section-phantom]{Descente et espaces algébriques}
\item \hyperref[spaces-perfect-section-phantom]{Catégories dérivées d'espaces}
\item \hyperref[spaces-more-morphisms-section-phantom]{Compléments sur les morphismes d'espaces}
\item \hyperref[spaces-flat-section-phantom]{Platitude sur les espaces algébriques}
\item \hyperref[spaces-groupoids-section-phantom]{Groupoïdes dans les espaces algébriques}
\item \hyperref[spaces-more-groupoids-section-phantom]{Compléments sur les groupoïdes dans les espaces}
\item \hyperref[bootstrap-section-phantom]{Amorçage}
\item \hyperref[spaces-pushouts-section-phantom]{Sommes amalgamées d'espaces algébriques}
\end{enumerate}
Thèmes de géométrie
\begin{enumerate}
\setcounter{enumi}{81}
\item \hyperref[spaces-chow-section-phantom]{Groupes de Chow des espaces}
\item \hyperref[groupoids-quotients-section-phantom]{Quotients de groupoïdes}
\item \hyperref[spaces-more-cohomology-section-phantom]{Compléments sur la cohomologie des espaces}
\item \hyperref[spaces-simplicial-section-phantom]{Espaces simpliciaux}
\item \hyperref[spaces-duality-section-phantom]{Dualité pour les espaces}
\item \hyperref[formal-spaces-section-phantom]{Espaces algébriques formels}
\item \hyperref[restricted-section-phantom]{Algébrisation des espaces formels}
\item \hyperref[spaces-resolve-section-phantom]{Résolution des surfaces revisitée}
\end{enumerate}
Théorie des déformations
\begin{enumerate}
\setcounter{enumi}{89}
\item \hyperref[formal-defos-section-phantom]{Théorie formelle des déformations}
\item \hyperref[defos-section-phantom]{Théorie des déformations}
\item \hyperref[cotangent-section-phantom]{Le complexe cotangent}
\item \hyperref[examples-defos-section-phantom]{Problèmes de déformation}
\end{enumerate}
Champs algébriques
\begin{enumerate}
\setcounter{enumi}{93}
\item \hyperref[algebraic-section-phantom]{Champs algébriques}
\item \hyperref[examples-stacks-section-phantom]{Exemples de champs}
\item \hyperref[stacks-sheaves-section-phantom]{Faisceaux sur les champs algébriques}
\item \hyperref[criteria-section-phantom]{Critères de représentabilité}
\item \hyperref[artin-section-phantom]{Axiomes d'Artin}
\item \hyperref[quot-section-phantom]{Espaces de Quot et de Hilbert}
\item \hyperref[stacks-properties-section-phantom]{Propriétés des champs algébriques}
\item \hyperref[stacks-morphisms-section-phantom]{Morphismes de champs algébriques}
\item \hyperref[stacks-limits-section-phantom]{Limites de champs algébriques}
\item \hyperref[stacks-cohomology-section-phantom]{Cohomologie des champs algébriques}
\item \hyperref[stacks-perfect-section-phantom]{Catégories dérivées de champs}
\item \hyperref[stacks-introduction-section-phantom]{Introduction aux champs algébriques}
\item \hyperref[stacks-more-morphisms-section-phantom]{Compléments sur les morphismes de champs}
\item \hyperref[stacks-geometry-section-phantom]{La géométrie des champs}
\end{enumerate}
Thèmes de théorie des espaces de modules
\begin{enumerate}
\setcounter{enumi}{107}
\item \hyperref[moduli-section-phantom]{Champs de modules}
\item \hyperref[moduli-curves-section-phantom]{Espaces de modules de courbes}
\end{enumerate}
Divers
\begin{enumerate}
\setcounter{enumi}{109}
\item \hyperref[examples-section-phantom]{Exemples}
\item \hyperref[exercises-section-phantom]{Exercices}
\item \hyperref[guide-section-phantom]{Guide bibliographique}
\item \hyperref[desirables-section-phantom]{Desiderata}
\item \hyperref[coding-section-phantom]{Style de programmation}
\item \hyperref[obsolete-section-phantom]{Obsolète}
\item \hyperref[fdl-section-phantom]{Licence de documentation libre GNU}
\item \hyperref[index-section-phantom]{Index généré automatiquement}
\end{enumerate}
\end{multicols}


\bibliography{my}
\bibliographystyle{amsalpha}

\end{document}
